\usepackage[T1]{fontenc}
\usepackage[utf8]{inputenc}
\usepackage{lmodern}
% La vía moderna de babel permite importar los rótulos españoles también
% en instalaciones que no incluyen el archivo tradicional spanish.ldf.
\IfFileExists{spanish.ldf}{%
  \usepackage[spanish,es-nodecimaldot,es-noquoting]{babel}%
}{%
  \usepackage{babel}
  \babelprovide[import,main]{spanish}%
}
\usepackage{amsmath,amssymb,amsthm}
\usepackage{mathtools}
\usepackage[margin=25mm,headheight=15pt,headsep=10mm,footskip=12mm]{geometry}
\usepackage{microtype}
\usepackage{xcolor}
\usepackage{booktabs,tabularx}
\usepackage{enumitem}
\usepackage{fancyhdr}
\usepackage{hyperref}
\usepackage{bookmark}

\definecolor{azulnota}{HTML}{203B55}
\definecolor{grisnota}{HTML}{59636D}
\hypersetup{
  colorlinks=true,
  linkcolor=azulnota,
  citecolor=azulnota,
  urlcolor=azulnota,
  pdftitle={Hipótesis de Riemann — Notas de estudio},
  pdfauthor={Alevi Peña},
  pdfsubject={Lectura de dos exposiciones del argumento de Alpöge para las funciones L de Dirichlet},
  pdfkeywords={Riemann, Dirichlet, Mellin, retículos, notas de estudio}
}
\pagestyle{fancy}
\fancyhf{}
\fancyhead[L]{\small\textcolor{grisnota}{Hipótesis de Riemann}}
\fancyhead[R]{\small\textcolor{grisnota}{Notas de estudio}}
\fancyfoot[C]{\small\thepage}
\renewcommand{\headrulewidth}{0.3pt}
\setlength{\parindent}{0pt}
\setlength{\parskip}{0.55em}
\setlength{\emergencystretch}{2em}
\setlist{topsep=0.4em,itemsep=0.3em,parsep=0pt}
\setlist[itemize,enumerate]{leftmargin=*}
\setlist[description]{style=nextline,leftmargin=1.4em,labelindent=0pt}
\newcolumntype{Y}{>{\raggedright\arraybackslash}X}
\setcounter{tocdepth}{2}
\numberwithin{equation}{section}
\allowdisplaybreaks[1]
\widowpenalty=10000
\clubpenalty=10000

\theoremstyle{plain}
\newtheorem{proposicion}{Proposición}[section]
\newtheorem{teorema}[proposicion]{Teorema}
\newtheorem{lema}[proposicion]{Lema}
\theoremstyle{definition}
\newtheorem{definicion}[proposicion]{Definición}
\newtheorem{ejercicio}{Ejercicio}[section]
\theoremstyle{remark}
\newtheorem*{observacion}{Observación}
\renewcommand{\proofname}{Demostración}
\renewcommand{\contentsname}{Índice}
\renewcommand{\refname}{Fuentes de esta entrega}

\newcommand{\N}{\mathbb{N}}
\newcommand{\Z}{\mathbb{Z}}
\newcommand{\Q}{\mathbb{Q}}
\newcommand{\R}{\mathbb{R}}
\newcommand{\C}{\mathbb{C}}
\newcommand{\OK}{\mathcal{O}_K}
\newcommand{\dd}{\,\mathrm{d}}
\DeclareMathOperator{\Ree}{Re}
\DeclareMathOperator{\Res}{Res}
\DeclareMathOperator{\Hom}{Hom}
\DeclareMathOperator{\length}{length}
\DeclareMathOperator{\ord}{ord}
\newcommand{\fuente}[1]{\par\small\textcolor{grisnota}{\textbf{Lectura de referencia.} #1}\par\normalsize}
